% ======================================================================
%  Section 4.7.1 --- Nominal case.  Corrected.
%
%  This is much the strongest section of the chapter and most of it needs
%  nothing. Nine items, of which four are errors, three are consistency
%  repairs following from edits we made upstream, and two are additions
%  that the numbers in your own table already support.
%
%  ERRORS
%  A. THE SECTION REFERS TO ITSELF. "The decisive comparison of the chapter
%     is accordingly the one made in Section~\ref{sec:ablation-results}."
%     That label belongs to 4.7, the section this subsection is inside, so
%     the sentence directs the reader to its own parent. It means 4.7.2.
%
%  B. THE THREE FIGURE CAPTIONS CONTRADICT THE TEXT. All three say
%     "feature-error norm $\lVert\mathbf{e}\rVert$". The text says, in
%     terms, "The quantity reported is the sum of squares
%     $\norm{\mathbf{e}}^{2}$ rather than the norm". The plots are of
%     ||e||^2. Captions corrected.
%
%  C. "no curve exhibits a plateau, a rebound or a secondary minimum",
%     followed two sentences later by "the feature error does not decay to
%     zero but settles on a floor". A floor is a plateau. The claim that
%     is meant, and that the curves support, is that none plateaus BEFORE
%     converging.
%
%  D. "the three variants ... Variant 3 recovers part of that loss". True
%     in rotation, FALSE in translation, where the ordering is
%     0.336 < 0.806 < 0.999 and Variant 3 is the worst of the three. The
%     paragraph discusses rotation only and does not say so. An examiner
%     reading the table will see it immediately, and silence on the point
%     reads worse than the result itself does. See (H): the translation
%     numbers have an explanation, and giving it turns the problem into a
%     finding.
%
%  CONSEQUENT REPAIRS (from edits made upstream in this pass)
%  E. THREE STALE LABELS.
%       sec:ablation-measure  -> sec:hermite-weight-limitation  (4.4.4;
%            we settled on YOUR label and dropped my two)
%       sec:implementation    -> eq:hessian-floor               (4.2.2;
%            the floor is an equation there, not a section elsewhere)
%       sec:ablation-protocol -> sec:metrics                    (4.6.4;
%            the second of the two uses, "the criterion stated in")
%
%  F. "decisive comparison" was removed from 4.6.1 for the reason that
%     4.7.2 reports Variants 2 and 3 differing by 7.6 per cent in rotation
%     and 7.2 per cent in translation IN OPPOSITE DIRECTIONS -- a tie. The
%     phrase survived here. Same treatment, or the chapter promises twice
%     what it delivers once.
%
%  G. THE "FEATURE ERROR IS NOT A PROXY" ARGUMENT IS NOW MADE TWICE. My
%     correction to 4.6.4 added it as policy, where the metrics are
%     declared. Here it should be the evidence for that policy, not a
%     second statement of it. Trimmed accordingly -- the measurement stays,
%     the restatement of the convergence criterion goes.
%
%  ADDITIONS YOUR OWN TABLE SUPPORTS
%  H. THE DEGRADATION IS ALMOST ENTIRELY IN TWO COMPONENTS. Of the
%     translational error, the depth term accounts for 90, 98 and 99 per
%     cent of the norm in Variants 1, 2 and 3; of the rotational error,
%     omega_x accounts for the bulk in all three. Weighting does not
%     degrade the pose uniformly -- it degrades depth and the rotation
%     coupled to it. That is worth a sentence: it is consistent with the
%     rank discussion of 4.5, and it explains (D).
%
%  I. AN ARITHMETIC CHECK YOU SHOULD RUN -- see the note at the foot.
%     The quoted gain values and the quoted intensity errors do not appear
%     to be consistent with each other.
% ======================================================================

\section{Results and Discussion}
\label{sec:ablation-results}

\subsection{Nominal case: behaviour without occlusion}
\label{sec:nominal}

The three variants are first compared on the clean target, with no occlusion
present. Two questions are settled by this run. The first is what the weighted
formulations cost when there is nothing to reject:
Remark~\ref{rem:baseline} establishes algebraically that $\mathbf{D}=\mathbf{I}$
reduces~\eqref{eq:robust-lm-law} to the control law of
Chapter~\ref{ch:herpvs}, but it says nothing about $\mathbf{D}\neq\mathbf{I}$
acting on uncorrupted measurements, and that is what the nominal run measures.
The second is a matter of attribution: the nominal behaviour of each variant
must be known before any difference observed under occlusion can be ascribed to
the treatment of the occlusion rather than to the weighting altering the
control law independently of it.

\begin{figure}[H]
    \centering

    \begin{minipage}[b]{0.23\textwidth}
        \centering
        \includegraphics[width=\textwidth]{Figs/C/fig1_a.PNG}

        \vspace{1mm}
        \textbf{(a)}
    \end{minipage}
    \hfill
    \begin{minipage}[b]{0.23\textwidth}
        \centering
        \includegraphics[width=\textwidth]{Figs/D/b.PNG}

        \vspace{1mm}
        \textbf{(b)}
    \end{minipage}
    \hfill
    \begin{minipage}[b]{0.23\textwidth}
        \centering
        \includegraphics[width=\textwidth]{Figs/D/c.PNG}

        \vspace{1mm}
        \textbf{(c)}
    \end{minipage}
    \hfill
    \begin{minipage}[b]{0.23\textwidth}
        \centering
        \includegraphics[width=\textwidth]{Figs/D/d.PNG}

        \vspace{1mm}
        \textbf{(d)}
    \end{minipage}

   \caption{Nominal case, experimental setup: (a)~desired image
   $\mathbf{I}^{*}$; (b)~current image $\mathbf{I}$ at the initial pose;
   (c)~initial image error $\mathbf{I}-\mathbf{I}^{*}$; (d)~final image error
   at convergence, Variant~3.}
   %%% (d) must say which variant. There are three final images and one
   %%% panel. Change "Variant~3" if the panel is from another run.
\label{fig:nom_setup}

\end{figure}

\paragraph{Conditions:}
The two poses, the textures and the iteration cap are those of
Section~\ref{sec:test-configurations}: the reference image $\mathbf{I}^{*}$ is
acquired from the unoccluded target at the desired pose, the images acquired
during servoing are taken from the same unoccluded target, so that the residual
arises solely from the pose error and no measurement is corrupted, and the
target covers the whole image at both poses so that no measurement is lost
off-frame. All remaining quantities are those fixed in
Section~\ref{sec:ablation-conditions} and are identical across the three
variants. That the three runs began from the same image is confirmed by the
initial intensity error, which took the same value $4.06\times10^{8}$ in all
three.
%%% (G, partial) The poses, the pixel displacements and the derivation of
%%% P = 264000 were stated in 4.6.2 and 4.6.3 and are here a third time.
%%% Cut to cross-references. What is kept is the one thing specific to this
%%% run and not stated anywhere else -- the identical initial intensity
%%% error, which is a genuine control and worth its sentence.
%%%
%%% IF YOU DECIDE NOT TO USE MY DRAFT OF 4.6.3, restore the two sentences
%%% giving t_0, omega_0, t*, omega* and the 41/65 pixel displacement here,
%%% or the poses will appear nowhere in the chapter.

\begin{figure}[htbp]
    \centering

    \begin{minipage}[b]{0.30\textwidth}
        \centering
        \includegraphics[width=\textwidth]{Figs/D/camera velocitiesV1.PNG}

        \vspace{1mm}
        \textbf{(a)}
    \end{minipage}
    \hfill
    \begin{minipage}[b]{0.30\textwidth}
        \centering
        \includegraphics[width=\textwidth]{Figs/D/poitiongerrorV1.PNG}

        \vspace{1mm}
        \textbf{(b)}
    \end{minipage}
    \hfill
    \begin{minipage}[b]{0.30\textwidth}
        \centering
        \includegraphics[width=\textwidth]{Figs/D/errV1.PNG}

        \vspace{1mm}
        \textbf{(c)}
    \end{minipage}

   \caption{Nominal case, Variant~1 (no weighting, $\mathbf{D}=\mathbf{I}$):
   (a)~camera velocities; (b)~translational and rotational positioning errors;
   (c)~feature error $\lVert\mathbf{e}\rVert^{2}$.}
   %%% (B) corrected here and in the two captions below.

   \label{fig:nom_v1}

\end{figure}

\begin{figure}[htbp]
    \centering

    \begin{minipage}[b]{0.30\textwidth}
        \centering
        \includegraphics[width=\textwidth]{Figs/D/cameraveloctiesV2PNG.PNG}

        \vspace{1mm}
        \textbf{(a)}
    \end{minipage}
    \hfill
    \begin{minipage}[b]{0.30\textwidth}
        \centering
        \includegraphics[width=\textwidth]{Figs/D/posiitionerrorV2.PNG}

        \vspace{1mm}
        \textbf{(b)}
    \end{minipage}
    \hfill
    \begin{minipage}[b]{0.30\textwidth}
        \centering
        \includegraphics[width=\textwidth]{Figs/D/ERRV2.PNG}

        \vspace{1mm}
        \textbf{(c)}
    \end{minipage}

    \caption{Nominal case, Variant~2 (Tukey weighting): (a)~camera velocities;
    (b)~translational and rotational positioning errors; (c)~feature error
    $\lVert\mathbf{e}\rVert^{2}$.}

    \label{fig:nom_v2}

\end{figure}

\begin{figure}[htbp]
    \centering

    \begin{minipage}[b]{0.30\textwidth}
        \centering
        \includegraphics[width=\textwidth]{Figs/D/caeravelocitiesV3.PNG}

        \vspace{1mm}
        \textbf{(a)}
    \end{minipage}
    \hfill
    \begin{minipage}[b]{0.30\textwidth}
        \centering
        \includegraphics[width=\textwidth]{Figs/D/poositionerrorV3.PNG}

        \vspace{1mm}
        \textbf{(b)}
    \end{minipage}
    \hfill
    \begin{minipage}[b]{0.30\textwidth}
        \centering
        \includegraphics[width=\textwidth]{Figs/D/ERROEV3.PNG}

        \vspace{1mm}
        \textbf{(c)}
    \end{minipage}

    \caption{Nominal case, Variant~3 (Hermite-informed weighting): (a)~camera
    velocities; (b)~translational and rotational positioning errors;
    (c)~feature error $\lVert\mathbf{e}\rVert^{2}$.}

     \label{fig:nom_v3}
\end{figure}

\paragraph{Results:}
Figures~\ref{fig:nom_v1}--\ref{fig:nom_v3} show the camera velocities, the
positioning error and the feature error against iteration, and
Table~\ref{tab:ch4-nominal} collects the final pose errors. All three variants
position the camera to better than one millimetre and better than half a
degree, which is the accuracy expected of a direct scheme at
convergence~\cite{key?}, and the decay is monotone in every case: no curve
plateaus before converging, rebounds, or settles on a secondary minimum.
%%% (C) "no curve exhibits a plateau" -> "no curve plateaus BEFORE
%%% converging". As written it was contradicted by the next paragraph,
%%% which reports the terminal floor.
%%% ~\cite{key?}: the sub-millimetre figure is Collewet and Marchand's, and
%%% the claim needs their reference or it reads as an assertion. Insert the
%%% key you use for photometric visual servoing in Chapter 3.
Two observations require comment, as both bear on how the remaining sections
must be read.

First, the feature error does not decay to zero but settles on a floor:
$4.35\times10^{7}$, $1.04\times10^{8}$ and $6.48\times10^{7}$ for
Variants~1,~2 and~3, while the commanded velocity falls to between
$1.4\times10^{-4}$ and $3.7\times10^{-4}\,\mathrm{m\,s^{-1}}$ and the pose
error becomes sub-millimetre. The quantity reported is the sum of squares
$\norm{\mathbf{e}}^{2}$ rather than the norm, so over $264\,000$ components
these floors correspond to root-mean-square residuals of $13$, $20$ and $16$
per component respectively, against $1149$ at the first iteration.
%%% Was "about 16 per component", which is Variant 3's value presented as
%%% though it were common to all three. sqrt(4.35e7/264000) = 12.8 and
%%% sqrt(1.04e8/264000) = 19.8, so the three differ by half again. Giving
%%% all three costs four words and removes a number an examiner can falsify
%%% with a calculator.
The run terminates not because the residual has been minimised but because the
adaptive gain has decayed with it: by~\eqref{eq:lambda} the gain is
proportional to the intensity error, $\lambda=\norm{\mathbf{e}_I}^{2}/N_{e}^{0.8}$,
so it falls from $\lambda\simeq53$ at the first iteration to $\lambda=0.041$ at
iteration~600. Since the commanded velocity is in addition proportional to the
residual itself, $\norm{\mathbf{V}_c}$ decays much faster than the error and
vanishes while a small misalignment remains. The minimisation is not at fault:
over the same run the feature error falls by three and a half orders of
magnitude, the commanded velocity from $0.566$ to
$2.3\times10^{-4}\,\mathrm{m\,s^{-1}}$, and the intensity error from $72.7$ to
$2.5$ grey levels root-mean-square. It is also worth recording that the damping
becomes inactive very early: by~\eqref{eq:mu}, $\mu\propto\norm{\mathbf{e}_I}^{10}$,
so $\mu$ falls from $1.9\times10^{-2}$ to $5.3\times10^{-18}$ over the same run.
The scheme is a Levenberg--Marquardt method only during the first few iterations
and a Gauss--Newton method thereafter, which is why the fixed floor added to the
damped Hessian in~\eqref{eq:hessian-floor} is what actually guarantees
invertibility once robust rejection is active.
%%% (E) sec:implementation -> eq:hessian-floor.

Second, and more consequentially, the feature error cannot be read as a proxy
for positioning accuracy across variants, which is why
Section~\ref{sec:metrics} admits it as a convergence indicator only. The
nominal run shows why directly. An alternative structure measure, reported in
Section~\ref{sec:hermite-weight-limitation}, attains a feature-error floor of
$6.59\times10^{7}$ against $6.48\times10^{7}$ here --- a difference of two per
cent --- yet its rotational pose error is smaller by a factor of four,
$0.088^{\circ}$ against $0.385^{\circ}$. The two rankings are not merely
uncorrelated but reversed: the variant with the larger residual is the more
accurately positioned of the two. The reason is that the residual of a weighted
variant is shaped not only by the misalignment but by which measurements the
estimator retains, so the floor mixes an accuracy term with a coverage term and
a variant can lower it by rejecting the measurements that are hardest to
satisfy. The pose error of Table~\ref{tab:ch4-nominal} is accordingly the
accuracy metric throughout this chapter.
%%% (E) sec:ablation-measure -> sec:hermite-weight-limitation.
%%% (G) The restatement of the convergence criterion and of the 10^4
%%% threshold is cut: 4.6.2 states it and 4.6.4 now carries the policy. What
%%% is kept here is the evidence, which belongs in the results.
%%% The direction of the discrepancy is now given. "differs by a factor of
%%% four" concealed that the discarded variant is the better positioned,
%%% which is the fact that makes the point. Concealing it would not have
%%% survived 4.4.4, where the same two numbers appear, and stating it here
%%% is stronger: a reversal is a much better argument than a decorrelation.

\paragraph{What the nominal case does and does not establish:}
The three variants do not coincide on a clean scene. Weighting is therefore not
transparent in the absence of outliers, and the differences in
Table~\ref{tab:ch4-nominal} are real rather than numerical. In rotation,
Variant~2 reaches $0.60^{\circ}$ against the baseline's $0.22^{\circ}$, which
is the cost anticipated in Section~\ref{sec:tukey-limitation}: while the pose
error is still large, valid measurements in photometrically sensitive regions
produce legitimately large residuals, a fraction of which a redescending
estimator rejects, so a weighted variant discards information the unweighted
scheme retains. Variant~3 recovers part of that loss, reaching $0.38^{\circ}$,
which is consistent with the structure normalisation restoring measurements
that the unnormalised estimator rejects on the strength of their residual
alone. In translation the ordering is the reverse: $0.336$, $0.806$ and
$0.999\,\mathrm{mm}$, with Variant~3 the least accurate of the three.
%%% (D) The reversal, stated. It is visible in the table on the facing page
%%% and passing over it in silence costs more than reporting it.

%%% (H) The explanation, which your own table supplies.
The two orderings are less contradictory than they appear once the error is
resolved into components. The degradation is not spread over the six degrees of
freedom but concentrated in two. Of the translational error, the depth term
$\Delta T_z$ accounts for $90$, $98$ and $99$ per cent of the norm in
Variants~1,~2 and~3 respectively; of the rotational error, $\Delta\omega_x$
accounts for the greater part in all three. Depth and the rotation about the
axis coupled to it are the directions least well constrained by a
fronto-parallel photometric measurement, since a small retreat and a small
tilt produce image changes of comparable form. Removing measurements can only
worsen the conditioning of directions that were already the weakest, and it is
in precisely those two that the weighted variants lose accuracy, while the four
well-constrained components remain at or below $0.2$ in every variant. The
penalty for weighting a clean scene is therefore not a uniform loss of
precision but a further softening of the two degrees of freedom that a direct
scheme constrains least securely. This is consistent with the rank argument of
Section~\ref{sec:robust-stability}, though it is not established by it: the
conditioning of $\mathbf{H}_{\mathbf{D}}$ is not measured here.
%%% The hedge in the last sentence is deliberate. 4.5 now says explicitly
%%% that the conditioning is not monitored, so this paragraph must not be
%%% allowed to imply that it was.

No ordering of the three is asserted on this evidence. These are single runs
from a single initial displacement, and differences of a few tenths of a degree
are of the same order as the spread a different displacement would plausibly
produce; establishing a ranking would require a study over several initial
poses, which is not undertaken here. What the nominal
case establishes is narrower and sufficient for what follows: that all three
variants converge, that none is destabilised by the weighting, that the penalty
for weighting a clean scene is confined to a few tenths of a degree in rotation
and a fraction of a millimetre in translation and falls almost entirely on two
degrees of freedom, and that this penalty must be subtracted from any advantage
observed under occlusion. The comparison on which the claim of the chapter
rests is accordingly the one made in Section~\ref{sec:occlusion}, where the
measurements are corrupted and the weighting has something to reject.
%%% (A) The self-reference is repaired: sec:ablation-results is THIS
%%% section, and the sentence means 4.7.2.
%%% (F) "The decisive comparison of the chapter" -> "The comparison on which
%%% the claim of the chapter rests", the same wording now used in 4.6.1. The
%%% phrase was removed there because 4.7.2 reports a tie between Variants 2
%%% and 3; it must not survive here.

\begin{table}[t]
    \centering
    \caption{Nominal case (unoccluded target): final pose error at
    iteration~600, relative to the desired pose $(0,0,1300)\,\mathrm{mm}$,
    $(0,0,0)^{\circ}$. Translations in millimetres, rotations in degrees.
    No run reached the absolute threshold on $\norm{\mathbf{e}}^{2}$; all three
    satisfied $\norm{\mathbf{V}_c}<10^{-3}\,\mathrm{m\,s^{-1}}$. Images of
    $320\times240$ pixels; feature dimension $264\,000$. One run per variant:
    see the discussion above on what may and may not be concluded.}
    \label{tab:ch4-nominal}
    \small
    \begin{tabular}{@{}lrrrrrrrr@{}}
        \toprule
        & \multicolumn{3}{c}{Translation (mm)}
        & \multicolumn{3}{c}{Rotation (deg)}
        & \multicolumn{2}{c}{Norm} \\
        \cmidrule(lr){2-4}\cmidrule(lr){5-7}\cmidrule(lr){8-9}
        Variant & $\Delta T_x$ & $\Delta T_y$ & $\Delta T_z$
                & $\Delta\omega_x$ & $\Delta\omega_y$ & $\Delta\omega_z$
                & $\norm{\Delta\mathbf{t}}$ & $\norm{\Delta\boldsymbol{\omega}}$ \\
        \midrule
        1 --- none    & $+0.053$ & $-0.140$ & $-0.301$ & $+0.168$ & $-0.139$ & $+0.015$ & $0.336$ & $0.219$ \\
        2 --- Tukey   & $+0.004$ & $-0.165$ & $-0.789$ & $+0.568$ & $-0.195$ & $-0.007$ & $0.806$ & $0.601$ \\
        3 --- Hermite & $+0.101$ & $-0.084$ & $-0.990$ & $+0.383$ & $-0.030$ & $+0.019$ & $0.999$ & $0.385$ \\
        \bottomrule
    \end{tabular}
\end{table}

% ======================================================================
%  (I) ONE ARITHMETIC CHECK BEFORE YOU SUBMIT. I could not make two of your
%  numbers agree, and an examiner who tries will hit the same wall.
%
%  The text states lambda = ||e_I||^2 / N_e^0.8 with N_e fixed over a run.
%  Then lambda is proportional to ||e_I||^2 and the two must fall by the
%  same factor. But:
%      lambda:  53 -> 0.041            a factor of 1293
%      ||e_I||: 72.7 -> 2.5 grey RMS   a factor of 29.1 in the norm,
%                                      hence 846 in the square
%  1293 against 846. Three possible explanations, and you can tell which in
%  one look at the log:
%    (i)  the 72.7 and 2.5 figures are from a different variant than the
%         lambda figures, in which case say which variant each belongs to;
%    (ii) N_e is the count of RETAINED measurements and varies over the run,
%         in which case lambda is not proportional to ||e_I||^2 and the
%         sentence "the gain is proportional to the intensity error" needs
%         qualifying;
%    (iii) one of the four numbers is a transcription slip.
%  Note that mu passes the same test: 29.1^10 = 4e14 against the quoted
%  1.9e-2/5.3e-18 = 3.6e15, which agrees to within the rounding of a
%  two-significant-figure quotation. So the mu sentence is sound and it is
%  specifically the lambda pair that needs checking.
%
%  WHAT I VERIFIED AND FOUND CORRECT
%  - All eight norms in Table 4.x, recomputed from their components. Every
%    one agrees to the last digit printed.
%  - The initial intensity error as a consistency check on the RMS:
%    sqrt(4.06e8 / 76800) = 72.7 grey levels, exactly the figure quoted in
%    the text. That is a good sign and worth keeping, since it ties two
%    independently reported numbers together.
%  - "more than three orders of magnitude below the floor" for the 10^4
%    threshold: 4.35e7/1e4 = 4350, so 3.6 orders. Correct.
%  - "falls by three and a half orders of magnitude": 1149^2 x 264000
%    = 3.5e11 down to 6.5e7 is a factor of 5400, i.e. 3.7 orders. Correct.
%
%  LABELS USED ABOVE
%    sec:test-configurations       4.6.3
%    sec:ablation-conditions       4.6.2
%    sec:metrics                   4.6.4
%    sec:tukey-limitation          4.3.3
%    sec:hermite-weight-limitation 4.4.4
%    sec:robust-stability          4.5
%    sec:occlusion                 4.7.2
%    sec:future                    future work
%    eq:hessian-floor              the floor added to 4.2.2
%
%  AND THE ONE LABEL STILL TO SETTLE: my correction to 4.5 refers to this
%  section as sec:results. You have labelled it sec:ablation-results. Change
%  the reference in 4.5 -- yours is the one in the thesis.
%
%  \begin{figure}[H] needs \usepackage{float} in the preamble.
% ======================================================================
